
%(BEGIN_QUESTION)
% Copyright 2011, Tony R. Kuphaldt, released under the Creative Commons Attribution License (v 1.0)
% This means you may do almost anything with this work of mine, so long as you give me proper credit

Les utvalgte deler av National Transportation Safety Board sin sikkerhetsstudie, \textit{Supervisory Control and Data Acquisition (SCADA) in Liquid Pipelines} (Document NTSB/SS-05/02 ; PB2005-917005), og svar på følgende spørsmål:

\vskip 10pt

Spørsmål \#22 i Safety Study Appendix E lister hovedfunksjoner i SCADA-systemer for rørledning.  Identifiser de vanligste funksjonene.  Spørsmål \#38 lister antall datapunkter i hvert svarendes SCADA-system.  Hvor datarike er noen av disse systemene?

\vskip 10pt

Spørsmål \#42 i Appendix E lister kommunikasjonsmedier brukt mellom RTU- og MTU-punkter.  Forklar hva hvert begrep betyr.  Hvor mange SCADA-systemer bruker ikke redundant (backup) kommunikasjon mellom RTU og MTU (hint: se spørsmål \#43)?

\vskip 10pt

Spørsmål \#19 og \#21 i Appendix E lister begrunnelser respondentene gir for hvorfor de vurderer oppgradering (eller er i gang med oppgradering) av SCADA-systemet sitt.  Hvilke begrunnelser oppgis?  Overrasker noen av dem deg?

\vskip 10pt

Seksjonen ``SCADA Screens and Graphics'' i kapittel 4 viser flere eksempler på grafiske skjermer i bruk for kontroll av væskerørledninger.  Hvor mye blankplass bør en skjerm ha for å unngå ``clutter''?  Hvordan bør farger velges for maksimal effekt?

\vskip 10pt

Resultatene i Appendix E lister en rekke leverandører (``vendors'') for SCADA-systemer brukt i rørledningskontroll (se spørsmål \#15).  Basert på leverandørnavn og antall installasjoner, hva er inntrykket av SCADA for rørledninger i USA: mye standardisering, eller stor systemvariasjon?  Er det en tydelig markedsleder i undersøkelsen?


\vskip 20pt \vbox{\hrule \hbox{\strut \vrule{} {\bf Forslag til sokratisk diskusjon} \vrule} \hrule}

\begin{itemize}
\item{} En interessant del av rapporten er ``Alarm Philosophy''.  Her beskrives hvordan alarmtilstander presenteres for operatører via SCADA, og utfordringer rundt håndtering av denne informasjonen.  Les og diskuter funnene.
\item{} En annen interessant del er ``Training and Selection'', der rapporten beskriver hvordan operatører trenes.  Det er interessant å sammenligne ulike treningsmetoder og modaliteter.  Les og diskuter funnene.
\item{} Det finnes en typografisk feil i bildeteksten til figur 4.10 -- finner du den?
\end{itemize}

\underbar{file i00277}
%(END_QUESTION)





%(BEGIN_ANSWER)


%(END_ANSWER)





%(BEGIN_NOTES)

I svarene på spørsmål \#22 ser vi funksjoner som:

\begin{itemize}
\item{} Fjernstyring av ventiler
\item{} Datainnsamling
\item{} Trendkurver
\item{} Lekkasjedeteksjon i rørledning
\end{itemize}

\vskip 10pt

I svarene på spørsmål \#38 ser vi at de fleste rørlednings-SCADA-systemer har {\it 5000 eller flere} telemetripunkter.

\vskip 10pt

En RTU er en Remote Terminal Unit, altså feltmontert SCADA-node der sensorer og ventilkontroller kobles inn.  En MTU står sentralt og henter data fra flere RTU-er.  Mest brukt kommunikasjonsform i denne undersøkelsen var leased lines (telefonlinjer), med satellitt og toveis radio på andre- og tredjeplass.  Omtrent 25\% (nøyaktig 25.9\%) av systemene har ingen nettverksredundans.  Blant de som har redundans, er ``dial-up'' (telefonmodem) mest brukt.  At prosentene summerer til over 100\% skyldes at samme system ofte bruker flere kommunikasjonsmetoder.  Spørsmål \#46 er også interessant: 15\% av systemene har opplevd total kommunikasjonssvikt på et tidspunkt.

\vskip 10pt

Maskinvare-obsolesens er den dominerende årsaken til SCADA-oppgradering i undersøkelsen.  Bedre sikkerhet rangerte også høyt blant respondentene som vurderte oppgradering.  Bare omtrent halvparten oppga at de faktisk byttet eller oppgraderte SCADA-systemet.

\vskip 10pt

En FAA-standard krever minst 40\% blankplass per skjerm for å unngå rot.  ISA sin mer liberale standard anbefaler 25\% til 40\% blankplass.  Farger bør brukes konsekvent (f.eks. ikke bruk rød både for alarmer og vanlige etiketter), og med tilstrekkelig kontrast (f.eks. unngå gul tekst på lys bakgrunn som i figur 4.7, eller lyseblå tekst på grå bakgrunn som i figur 4.8).  ISA anbefaler 5 til 7 farger.  En interessant detalj i seksjonen var at fargeblinde operatører hadde problemer med å skille enkelte farger, og at flere av selskapene ikke testet for fargeblindhet.

\vskip 10pt

Stor variasjon i SCADA-navn vises i denne undersøkelsen: 41 ulike av 98 respondenter.  Antall installasjoner per navn er relativt lavt, noe som tyder på lite standardisering i bransjen.  Mest utbredt (13 installasjoner) er Valmet Automation.  Wonderware og Telvent er nest mest utbredt (8 installasjoner hver).







\vskip 20pt \vbox{\hrule \hbox{\strut \vrule{} {\bf Forslag til sokratisk diskusjon} \vrule} \hrule}

\begin{itemize}
\item{} Gi eksempler på dårlige fargevalg i SCADA-skjermer.
\item{} Gi eksempler på ``clutter'' i SCADA-skjermer.
\item{} Forklar hvorfor du tror så mange SCADA-design er svake.  Hvorfor så mange dårlige fargevalg og rotete skjermer?
\end{itemize}







\vfil \eject

\noindent
{\bf Forberedelsesquiz:}

Basert på NTSB-undersøkelsen av SCADA-systemer brukt i kontroll av væskerørledninger i USA (2002), hvor standardiserte er disse SCADA-systemene?

\begin{itemize}
\item{} Høyt standardisert; de fleste rørledninger styres av Allen-Bradley SCADA-systemer
\vskip 5pt 
\item{} Høyt standardisert; de fleste rørledninger styres av Siemens SCADA-systemer
\vskip 5pt 
\item{} Nøyaktig ti ulike produsenter representert i undersøkelsen
\vskip 5pt 
\item{} Lite standardisert; stor variasjon av produsentløsninger representert
\vskip 5pt 
\item{} Ingen standardisering; ingen produsent var representert mer enn én gang
\end{itemize}


%INDEX% Reading assignment: NTSB Report, SCADA in Liquid Pipelines

%(END_NOTES)
